\section{O1：Scaling Law的新范式}

O1的发布是本次圆桌讨论的核心话题。三位嘉宾从不同角度给出了高度一致的积极评价。

\subsection{System 1 + System 2的统一}

\begin{importantbox}{蒋大昕的分析}
O1的核心意义有二：
\begin{enumerate}
\item 第一次证明语言模型可以拥有System 2能力——不再是直线性思维（System 1），而是能探索不同路径、自我反思、自我纠错，直到找到正确答案
\item 将模仿学习（预训练）与强化学习结合，使模型同时具备System 1和System 2能力
\end{enumerate}
\end{importantbox}

\begin{knowledgebox}{System 1 vs. System 2}
\textbf{System 1}（快思考）：GPT-4的思维方式——即使拆解复杂问题为多步，仍是直线性的next-token prediction。

\textbf{System 2}（慢思考）：O1的新能力——可以探索多条路径、回溯、纠错，直到找到正确方案。这是归纳世界、解决新问题的基础能力。
\end{knowledgebox}

\subsection{RL Scaling：新的Scaling维度}

蒋大昕进一步指出，O1带来了一个\textbf{RL Scaling}的新范式：

\begin{itemize}
\item 过去DeepMind的强化学习（AlphaGo/AlphaFold/AlphaGeometry）都是为特定场景设计
\item O1首次在通用性和泛化性上实现了大规模RL Scaling
\item 它还只是开端（preview），但已经展示了一条上限很高的道路
\end{itemize}

\subsection{AGI的等级跃迁}

\begin{knowledgebox}{朱军的AGI L1--L5分级}
\begin{itemize}
\item \textbf{L1 聊天机器人}：ChatGPT时代
\item \textbf{L2 推理者}：O1时代——在特定（Narrow）场景下已实现L2
\item \textbf{L3 智能体}：从数字世界走向物理世界，去交互、去改变
\item \textbf{L4 创新者}：发现新知识、创造新事物
\item \textbf{L5 组织者}：协同组织、高效运转
\end{itemize}
每一级都有从Narrow到General的跃迁。O1代表着从L1向L2的显著质变。
\end{knowledgebox}

\subsection{解决Scaling Wall}

杨植麟指出O1提升了AI的\textbf{上限}：

\begin{importantbox}{从5\%到10倍}
O1之前，大家担心数据枯竭（数据墙）导致Scaling Law失效。O1证明了通过强化学习可以继续Scaling——从``提升5--10\%的生产力''的上限提升到``10倍GDP''的可能性。AI七八十年历史上唯一有效的就是Scaling，O1为Scaling指出了新的维度。
\end{importantbox}

\subsection{本章小结}

O1的意义：(1) 语言模型获得System 2能力；(2) RL Scaling新范式打破数据墙；(3) AGI从L1迈向L2。


